%! ~~~ Packages Setup ~~~ 
\documentclass[]{../../../../tex/classes/styledArticle}
\usepackage{../../../../tex/packages/hebrewSupport}
\usepackage{../../../../tex/packages/mathShortcuts}
\usepackage{../../../../tex/packages/theoremsSupport}
\usepackage{../../../../tex/packages/enfancysections}

%! ~~~ Document ~~~

\author{שחר פרץ}
\title{\textit{לינארית 2א $\sim$ משפט הפירוק הספקטרלי להעתקות כלליות}}
\date{22 ביוני 2025}
\begin{document}
	\maketitle
	\textbf{מבוסס על הקלטה 18 של הקורס ב־2024-2025}
	
	\textbf{מרצה: ענת אמיר}
	
	המטרה: להבין לאילו העתקות בדיוק מתקיים המשפט הספקטרלי. מעל הממשיים, הבנו שאילו העתקות צמודות לעצמן. אז מה קורה מעל המרוכבים? 
	
	\theo{(משפט z). יהי $V$ ממ''פ סופי ויהי $\vphi \in V^*$. אז קיים ויחיד וקטור $u \in V$ שמקיים $\forall v \in V \co \vphi(v) = \mut{v}{u}$. }\begin{proof}
		\textbf{קיום. }יהי $B = (b_i)_{i = 1}^{n}$ בסיס אורתונורמלי של $V$ (הוכחנו קיום בהרצאות קודמות). נסמן $u = \sum_{i =1}^{n}\ol{\phi(b_i)} b_i$. בכדי להראות $\forall v \in V \co \phi(v) = \mut{v}{u}$ מספיק להראות תכונה זו לאברי הבסיס $B$, כלומר נראה ש־$\forall 1 \le j \le n \co \phi(b_j) = \mut{b_j}{u}$. ואכן: 
		\[ \mut{b_j}{u} = \mut{b_j}{\sumni \ol{\phi(b_i)}b_i} = \sum_{i = 1}^{n}\underbrace{\ol{\ol{\phi(b_i)}}}_{b_i}\underbrace{\mut{b_j}{b_i}}_{\dj_{ij}} = b_j \quad \top \]
		\textbf{יחידות: }אם קיים וקטור נוסף שעבורו $\forall v \in V \co \phi(v) = \mut{v}{w}$ אז בפרט עבור $v = u - w$ נקבל: 
		\[ \phi(v) = \mut{v}{w} = \mut{v}{u} \implies \mut{v}{u - w} = 0 \implies 0 = \mut{u - w}{u - w} = \norm{v - w}^{2} = 0 \implies v - w = 0 \implies v = w \]
		סה''כ הוכחנו קיום ויחידות. 
	\end{proof}
	
	\theo{יהי $V$ ממ''פ מנ''ס ותהי $T \co V \to V$ לינארית. אז קיימת ויחידה $T^* \co V \to V$ ומקיימת $\forall u ,v \in V \co \mut{Tu}{v} = \mut{u}{T^*v}$. }
	\defi{ההעתקה $T^*$ לעיל נקראת \textit{ההעתקה הצמודה ל־$T$}. }
	\begin{proof}
		לכל $v \in V$, נתבונן בפונקציונל הלינארי $\phi_V \in V^*$ המוגדר ע''י $\forall u \in V \co \phi_V(u) = \mut{Tu}{v}$. ממשפט ריס קיים ויחיד $T^*v \in V$ שעבורו $\forall u \in V \co \mut{Tu}{v} = \phi_V(u) = \mut{u}{T^*v}$. כלומר, ההעתקה $T^* \co V \to V$, קיימת ויחידה, ונותר להראות שהיא לינארית. עבור $v, w \in V$ ועבור $\ag, \bg \in \F$ מתקיים: 
		\[ \forall u \in V \co \mut{u}{T^* (\ag v + \bg w)} = \mut{Tu}{\ag v + \bg w} = \bar \ag \mut{Tu}{v} + \bar \bg \mut{Tu}{v} = \bar \ag \mut{u}{T^* v} + \bar \bg \mut{u}{T^*w} = \mut{u}{\ag T^*u + \bg T^*w} \]
		מסך נסיק ש־$T^*(\ag v + \bg w) = \ag T^* u + \bg T^* w$ מנימוקים דומים. 
	\end{proof}
	
	\textbf{דוגמאות. }
		מעל $\C^n$, עם המ''פ הסטנדרטי, נגדיר ט''ל $T_A \co \C^n \to \C^n$ עבור $A \in M_n(\C)$ מוגדרת ע''י $T_A(x) = Ax$. אז: 
		\[ \forall x, y \in \C^n \co \mut{T_A(x)}{y} = \mut{Ax}{y} = \ol{(Ax)^T} \cdot y = \ol{x^T}\cdot\ol{A^T}y\cdot = \ol{x^T}T_{\ol{A^T}}(y) = \mut{x}{T_{\ol{A^T}}y} \]
		כלומר, $(T_A)^* = T_{A^*}$ כאשר $A^* = \ol{A^T}$, וקראנו לה המטריצה הצמודה. 
	
	\textit{הערה: ההוכחה הזו הפוכה כי צריך לינאריות ברכיב הראשון ולא בשני והמרצה התבלבלה אבל אין לי כוח לתקן. עדיין $(T_A)^* = T_{A^*}$. }
	
	נבחין שהעתקה נקראת צמודה לעצמה אממ $T^* = T$. 
	
	עוד נבחין שעבור העתקה הסיבוב $T \co \R^2 \to \R^2$ בזווית $\tg$, מתקיים ש־$T^*$ היא הסיבוב ב־$-\tg$, וכן היא גם ההופכית לה. כלומר $(T_{\tg})^* = T_{-\tg} = (T_\tg)\op$. זו תכונה מאוד מועילה וגם נמציא לה שם במועד מאוחר יותר. 
	
	\theo{(תכונות ההעתקה הצמודה) יהי $V$ ממ''פ ותהיינה $T, S \co V \to V$ זוג העתקות לינאריות. נבחין ש־: 
	\begin{enumerate}[(A)]
		\item \hfil $(T^*)^* = T$
		\item \hfil $(T \circ S)^* = S^* \circ T^*$
		\item \hfil $(T + S)^* = T^* + S^*$
		\item \hfil $\forall \lg \in \F \co (\lg T)^* = \bar \lg (T^*)$
	\end{enumerate}}
	``זה אחד וחצי לינאריות''
	
	\begin{proof}\,
		\begin{enumerate}[A)]
			\item \hfil $\forall u, v \in V \co \mut{T^* u}{v} = \ol{\mut{v}{T^* u}}  = \ol{\mut{Tv}{u}}  = \mut{u}{Tv} \implies (T^*)^* = T$
			\item \hfil $\mut{(T \circ S) u}{v} = \mut{Su}{T^*v} = \mut{u}{S^* T^*} \implies (TS)^* = T^*S^*$
			\item \hfil $\mut{(T + S)u}{v} = \mut{Tu}{v} + \mut{Su}{v} = \mut{u}{T^*v} + \mut{u}{S^*v}  = \mut{u}{T^*v + S^*v}$
			\item כנ''ל
		\end{enumerate}
	\end{proof}
	
	\theo{יהי $V$ ממ''פ נ''ס ותהי $T \co V \to V$ לינאריות. אם $B = (b_i)_{i = 1}^{n}$ בסיס אורתוגונלי לו''ע של $T$, אז $\forall 1 \le i \le n$ ו''ע של ההעתקה הצמודה. }
	כלומר: אם מתקיים המשפט הספקטרלי, אז הבסיס שמלכסן אורתוגונלית את $T$ מלכסן אורתוגונלית את הצמודה. 
	\begin{proof}
		יהי $i \in [n]$ ונסמן בעבורו את $\lg_i$ הע''ע המתאים לו''ע $b_i$. עבור $i \neq j \in [n]$ נחשב את $\mut{b_i}{T^*b_j}$: 
		\[ \mut{b_i}{T^*b_j} = \ol{\mut{Tb_i}{b_j}} = \ol{\mut{\lg _i b_i}{b_j}} = \lg_i \mut{b_i}{b_j} = 0 \]
		לכן $T^*b_j \in (\Sp\{b_i\}_{i = 1}^{n})^{\perp} \seq \Sp\{b_j\}$. משיקולי ממדים, הפריסה מממד $n - 1$ ולכן המשלים האורתוגונלי שלו מממד $1$ ולכן השוויון. סה''כ $T^* b_j \in \Sp\{b_j\}$ ולכן $b_j$ ו''ע של $T^*$ כדרוש. 
	\end{proof}
	
	\textbf{מסקנה. }אאם $V$ ממ''פ נ''ס ו־$T \co V \to V$ ט''ל עם בסיס מלכסן אורתוגונלי, אז $T, T^*$ מתחלפות כלומר $TT^* = T^*T$. \begin{proof}
		לפי הטענה הקודמת כל $b_i$ הוא ו''ע משותף ל־$T$ ול־$T^*$, ולכן: 
		\[ TT^*(b_i) = T(T^*(b_i)) = \bg_i T^*(b_i) = \bg_i \dg_i b_i = \ag_i \bg_i b_i = \ag_i T^(b_i) = T^*T(b_i) \]
		העתקה מוגדרת לפי מה שהיא עושה לבסיס ולכן $TT^* = T^*T$. 
	\end{proof}
	\defi{העתקה כזו המקיימת $AA^* = A^*A$ נקראת \textit{נורמלית}). }
	
	עתה, ננסה להראות שכל העתקה נורמלית מקיימת את המפשט הספקטרלי. 
	
	\theo{(המשפט הספקטרלי) יהי $V$ ממ''פ נוצר סופית מעל $\C$, ותהי $T \co V \to V$ לינארית. אז קיים בסיס אורתוגונלי של ו''ע של $T$ אמ''מ $T$ נורמלית. }
	\lem{יהי $V$ ממ''פ ותהיינה $S_1, S_2 \co V \to V$  זוג ט''ל צמודות ולעמן ומתחלפות (כלומר $S_1S_2 = S_2 S_1$). אז קיים בסיס אורתוגונלי של $V$ שמורכב מו''עים משופים ל־$S_1$ ול־$S_2$. }\begin{proof}
		ידוע ש־$S_1$ צמודה לעצמה, לכן לפי המשפט הספקטרלי להעתקות צמודות לעצמן (לא מעגלי כי הוכח בנפרד בהרצאה הקודמת), קיים לה לכסון אורותגונלי ובפרט $S_1$ לכסינה. נציג את $V$ כ־$V = \bigoplus_{i = 1}^{m} \ker(S_1 - \lg_iI)$, כאשר $\lg _1 \dots \lg_m$ הע''עים השונים של $S_1$. לכל $1 \le i \le m$  מתקיים ש־$V_{\lg_i}$ (המרחב העצמי) הוא $S_1$־אינווריאנטי שהרי אם $v \in V_{\lg_i}$ ונחשב: 
		\[ S_1(S_2 v) = S_2(S_1 v) = S_2(\lg_i v) = \lg _i S_2v \implies S_2 v \in V_{\lg_i} \]
		כאשר $S_2|_{V_{\lg _i}} \co V_{\lg_i} \to V_{\lg_i}$ צמודה לעצמה, ולכן המפשט הספקטרלי לצמודות לעצמן אומר שבתוך $V_{\lg_i}$ ישנו בסיס אורתוגונלי של ו''עים מ־$S_2$. האיחוד של כל הבסיסים הללו מכל מ''ע של $S_1$ יהיה בסיס אורתוגונלי של ו''עים משותפים ל־$S_1$ ול־$S_2$. 
	\end{proof}
	\begin{proof}[הוכחת המשפט הספקטרלי. ]\,
		\begin{itemize}
			\item[$\implies$] לפי המסקנה הקודמת, אם ישנו לכסון אורתוגונלי $T$ בהכרח נורמלית. 
			\item[$\impliedby$] נגדיר $S_1 = \frac{T + T^*}{2}, \ S_2 = \frac{T - T^*}{2i}$. הן וודאי צמודות לעצמן מהלינאריות וכל השטויות ממקודם, והן גם מתחלפות אם תטרחו להכפיל אותן. מהטענה קיים ל־$V$ בסיס אורתוגונלי של ו''עים משותפים ל־$S_1, S_2$ ונסמנו $\{b_i\}^{n}_{i = 1}$ וגם $S_1b_i = \ag_i b_i, \ S_2b_i= \bg_i b_i$. אפשר גם לטעון ש־$\ag_i, \bg_i \in \R$ אבל זה לא מועיל לנו. נשים לב ש־$T = S_1 + iS_2$, כלומר $\forall i \in [n] \co T(b_i) = S_1(b_i) + iS_2(b_i) = \ag_i b_i + i\bg_i b_i = (\ag + i\bg_i)b_i$ וזהו בסיס אורתוגונלי של ו''עים של $T$. 
		\end{itemize}
	\end{proof}
	
	למעשה, הבנו מהפירוק של $S_1, S_2$ ש־$S_1$ נותנת את החלק הממשי של הע''ע ו־$S_2$ את החלק המדומה. 
	
	``אגב – לא השתמשתי במשפט היסודי של האלגברה''
	
	
	
	
	\ndoc
\end{document}
